\chapter{\ifenglish Introduction\else บทนำ\fi}

\section{\ifenglish Project rationale\else ที่มาของโครงงาน\fi}
ในปัจจุบันองค์กรและสถาบันการศึกษาต่างๆ รวมถึงมหาวิทยาลัยเชียงใหม่มีการใช้พลังงานไฟฟ้า 
เพิ่มขึ้นอย่างต่อเนื่อง ส่งผลให้เกิดการปล่อยก๊าซคาร์บอนไดออกไซด์ในปริมาณมาก ถึงแม้ว่า 
มหาวิทยาลัยเชียงใหม่จะมีการติดตั้งระบบผลิตไฟฟ้าจากเซลล์แสงอาทิตย์ (Solar Cell) แล้วก็ตาม 
เนื่องจากทางมหาวิทยาลัยเชียงใหมยังไม่มีระบบจัดเก็บพลังงานสำรอง (Battery) 
ทำให้การผลิตไฟฟ้าไม่สอดคล้องกับการใช้งานจริงส่งผลให้มหาวิทยาลัยเชียงใหม่ต้องพึ่งพาพลังงานไฟฟ้าจาก
ภายนอก (PEA) มากกว่าการใช้พลังงานไฟฟ้าที่ผลิตได้ นอกจากนี้ การผลิตไฟฟ้าจากพลังงานแสงอาทิตย์ 
(Solar Cell) ยังมีความไม่แน่นอน เนื่องจากขึ้นอยู่กับสภาพอากาศที่เปลี่ยนแปลงลงตลอดเวลา
ทำให้การวางแผนการใช้พลังงานไฟฟ้าเป็นไปได้อย่างยากลำบาก 
ดังนั้นเพื่อให้สามารถลดการใช้พลังงานจากภายนอก (PEA) ได้ จำเป็นต้องมีการวางแผน 
และบริหารจัดการพลังงานให้มีประสิทธิภาพมากขึ้น หนึ่งในวิธีการแก้ปัญหาคือการใช้ข้อมูลสภาพอากาศ 
และปัจจัยที่มีผลต่อการผลิตไฟฟ้าจากเซลล์แสงอาทิตย์ (Solar Cell) มาวิเคราะห์และพยากรณ์ 
ปริมาณพลังงานที่สามารถผลิตได้ล่วงหน้า  
เพื่อให้สามารถพยากรณ์การผลิตพลังงานจากพลังงานแสงอาทิตย์ (Solar Cell) ได้อย่างแม่นยำขึ้น 
เราจึงมีการสร้างสถานีอุตุนิยมวิทยาขนาดย่อมขึ้นมาสำหรับวัดสภาพอากาศ โดยสถานีจะทำการเก็บข้อมูล 
เช่น ความเข้มแสง อุณหภูมิหน้าแผง ความชื้น ความเร็วลม และทิศทางลม โดยข้อมูลเหล่านี้จะถูก 
นำมาวิเคราะห์และใช้สำหรับพยากรณ์ปริมาณพลังงานไฟฟ้าที่จะผลิตได้ในแต่ละช่วงเวลา 
การใช้สถานีอุตุนิยมวิทยาขนาดย่อมเพื่อพยากรณ์ปริมาณการผลิตพลังงานจากเซลล์แสงอาทิตย์ 
(Solar Cell) จะช่วยให้มหาวิทยาลัยเชียงใหม่สามารถวางแผนการใช้พลังงานได้อย่างเหมาะสม 
ลดการพึ่งพาพลังงานไฟฟ้า จากภายนอก (PEA) และเพิ่มประสิทธิภาพการใช้พลังงานที่ผลิตได้ 
นอกจากนี้ยังช่วยลดการปล่อยแก๊สคาร์บอนไดออกไซด์ 

\section{\ifenglish Objectives\else วัตถุประสงค์ของโครงงาน\fi}
\begin{enumerate}
    \item ออกแบบและพัฒนาสถานีวัดสภาพอากาศขนาดเล็ก
    \item ออกแบบและพัฒนาการเก็บข้อมูลสภาพอากาศทีวัดได้จากสถานีวัดสภาพอากาศขนาดเล็ก
    \item อกแบบและพัฒนาระบบพยากรณ์ปริมาณการผลิตไฟฟ้าจากแผงโซลลาร์เซลล์ล่วงหน้า 24 ชั่วโมง
\end{enumerate}




\section{\ifenglish Project scope\else ขอบเขตของโครงงาน\fi}

\subsection{\ifenglish Hardware scope\else ขอบเขตด้านฮาร์ดแวร์\fi}
\begin{enumerate}
    \item การออกแบบและพัฒนาสถานีวัดสภาพอากาศขนาดย่อมสามารถวัดค่าอุณหภูมิ ความชื้น ความกดอากาศ ความเข้มแสง ความเร็วลม และทิศทางลม และอุณหภูมิของแผงโซลาร์เซลล์ได้
    \item การออกแบบและพัฒนาสถานีวัดสภาพอากาศขนาดย่อมสามารถส่งข้อมูลที่วัดได้ไปยังระบบจัดเก็บข้อมูลได้ทุกๆ 1 ชั่วโมง
\end{enumerate}

\subsection{\ifenglish Software scope\else ขอบเขตด้านซอฟต์แวร์\fi}
\begin{enumerate}
    \item ระบบสามารถจัดเก็บข้อมูลสภาพอากาศที่วัดได้จากสถานีวัดสภาพอากาศขนาดย่อมได้
    \item ระบบสามารถแสดงผลข้อมูลสภาพอากาศที่วัดได้จากสถานีวัดสภาพอากาศขนาดย่อมได้
    \item ระบบสามารถแสดงข้อมูลการพยากรณ์ปริมาณการผลิตไฟฟ้าจากแผงโซลาร์เซลล์ล่วงหน้า 24 ชั่วโมงได้
    \item ระบบสามารถตรวจสอบสถานะการทำงานของสถานีวัดสภาพอากาศขนาดย่อมได้
\end{enumerate}

\subsection{\ifenglish Model scope\else ขอบเขตด้าน Machine Learning Model\fi}
\begin{enumerate}
    \item การทำความสะอาดข้อมูล (Data Cleaning)
    \item การเลือกคุณลักษณะที่สำคัญต่อการพยากรณ์ปริมาณการผลิตไฟฟ้า (Feature Selection)
    \item โมเดลสามารถพยากรณ์ปริมาณการผลิตไฟฟ้าจากแผงโซลาร์เซลล์ล่วงหน้า 24 ชั่วโมงได้
\end{enumerate}




\section{\ifenglish Expected outcomes\else ประโยชน์ที่ได้รับ\fi}
\begin{itemize}
    \item[1.] ผู้ใช้สามารถวางแผนการใช้พลังงานไฟฟ้าได้อย่างมีประสิทธิภาพมากขึ้น โดยสามารถดูการพยากรณ์ปริมาณการผลิตไฟฟ้าจากแผงโซลาร์เซลล์ได้ล่วงหน้า 24 ชั่วโมง
    \item[2.] ช่วยลดการพึ่งพาพลังงานไฟฟ้าจากภายนอก (PEA) และเพิ่มประสิทธิภาพการใช้พลังงานที่ผลิตได้จากแผงโซลาร์เซลล์
    \item[3.] ช่วยลดการปล่อยแก๊สคาร์บอนไดออกไซด์ และส่งเสริมการใช้พลังงานสะอาดจากแผงโซลาร์เซลล์
\end{itemize}
    



\section{\ifenglish Technology and tools\else เทคโนโลยีและเครื่องมือที่ใช้\fi}

\subsection{\ifenglish Hardware technology\else เทคโนโลยีด้านฮาร์ดแวร์\fi}
\begin{itemize}
    \item[] \textbf{1.5.1 เซนเซอร์ (Sensor) }

    \begin{itemize}
    \item[] \textbf{1.5.1.1 เซนเซอร์วัดอุณหภูมิ ความชื้นและความกดอากาศ BME280 }
    \end{itemize}
    \begin{itemize}
    \item[] \textbf{1.5.1.2 เซนเซอร์วัดความเข้มแสง VEML7700 } 
    \end{itemize}
    \begin{itemize}
    \item[] \textbf{1.5.1.3 เซนเซอร์วัดความเร็วลม Wind Speed }
    \end{itemize}
    \begin{itemize}
    \item[] \textbf{1.5.1.4 เซนเซอร์วัดทิศทางลม Wind Direction }
    \end{itemize}
    \begin{itemize}
    \item[] \textbf{1.5.1.5 เซนเซอร์วัดอุณหภูมิ PT100 }
    \end{itemize}
    \begin{itemize}
    \item[]  \textbf{1.5.1.6 ไมโครคอนโทรลเลอร์ (Microcontroller) ESP32 }
    \end{itemize}
    \begin{itemize}
    \item[] \textbf{1.5.1.7 ซิมการ์ด (SIM Card) สำหรับการเชื่อมต่ออินเทอร์เน็ตผ่านเครือข่ายมือถือ SIM800L }
    \end{itemize}

\end{itemize}


\subsection{\ifenglish Software technology\else เทคโนโลยีด้านซอฟต์แวร์\fi}
\begin{itemize}
\item[] \textbf{1.5.2.1 Frontend }
    \begin{itemize}
        \item React: ใช้สำหรับสร้าง User Interface (UI)
        \item vite: เครื่องมือสำหรับรัน Development Server และ Build โปรเจค
        \item Figma: ใช้สำหรับออกแบบ UI และ UX ของระบบ
    \end{itemize}
\end{itemize}
\begin{itemize}
\item[] \textbf{1.5.2.2 Backend }
    \begin{itemize}
        \item FastAPI: ใช้สำหรับสร้าง RESTful API สำหรับจัดการข้อมูลสภาพอากาศ ข้อมูลการพยากรณ์ปริมาณการผลิตไฟฟ้าจากแผงโซลาร์เซลล์ และข้อมูลสถานะการทำงานของสถานีวัดสภาพอากาศขนาดย่อม
        \item PostgreSQL: ฐานข้อมูลสำหรับจัดเก็บข้อมูลสภาพอากาศและข้อมูลการพยากรณ์ปริมาณการผลิตไฟฟ้าจากแผงโซลาร์เซลล์
    \end{itemize}
\end{itemize}
\begin{itemize}
\item[] \textbf{1.5.2.3 Server  }
    \begin{itemize}
        \item Server ภาควิชาวิศวกรรมคอมพิวเตอร์ คณะวิศวกรรมศาสตร์ มหาวิทยาลัยเชียงใหม่ 
    \end{itemize}
\end{itemize}
\begin{itemize}
\item[] \textbf{1.5.2.4 Machine Learning Model }
    \begin{itemize}
        \item  Google Colab: ใช้สำหรับพัฒนาและทดสอบ Machine Learning Model สำหรับพยากรณ์ปริมาณการผลิตไฟฟ้าจากแผงโซลาร์เซลล์ล่วงหน้า 24 ชั่วโมง
    \end{itemize}
\end{itemize}





\section{\ifenglish Project plan\else แผนการดำเนินงาน\fi}

\begin{plan}{11}{2024}{3}{2026}
    \planitem{11}{2024}{11}{2024}{ปรึกษาอาจารย์ที่ปรึกษา}
    \planitem{11}{2024}{1}{2025}{เก็บข้อมูลและศึกษาข้อมูลเกี่ยวกับโครงงาน}
    \planitem{1}{2025}{1}{2025}{สรุปขอบเขตของโครงงาน}
    \planitem{12}{2025}{3}{2025}{ศึกษาอุปกรณ์ที่ใช้ในการพัฒนาสถานีและซอฟต์แวร์ที่จะนำมาพัฒนาระบบ}
    \planitem{3}{2025}{3}{2025}{ตรวจสอบสถานที่สำหรับการติดตั้งสถานี}
    \planitem{11}{2025}{2}{2026}{สร้างสถานีและออกแบบระบบจัดเก็บข้อมูล}
    \planitem{12}{2025}{2}{2026}{สร้าง Machine Learning Model}
    \planitem{2}{2026}{2}{2026}{ทดสอบการทำงานของสถานีและระบบจัดเก็บข้อมูล}
    \planitem{2}{2026}{2}{2026}{ทดสอบการทำงานของ Machine Learning Model}
    \planitem{3}{2026}{3}{2026}{ทดสอบการทำงานของระบบทั้งหมดและปรับปรุงแก้ไขตามผลการทดสอบ}
    \planitem{3}{2026}{3}{2026}{ทำโครงงานและเขียนรายงาน}
\end{plan}
หมายเหตุ: ในช่วงระยะเวลาเดือนเมษายน พ.ศ.2568 ถึง เดือนตุลาคม พ.ศ.2568 เป็นช่วงเวลาที่ไม่มีการดำเนินงานเนื่องจากสมาชิกทั้ง 3 คนได้ไปฝึกสหกิจศึกษา

\section{\ifenglish Roles and responsibilities\else บทบาทและความรับผิดชอบ\fi}
\begin{itemize}
    \item[1.] นางสาวธณิชา มังกร: รับผิดชอบในการออกแบบและพัฒนาระบบจัดเก็บข้อมูลสภาพอากาศที่วัดได้จากสถานีวัดสภาพอากาศขนาดย่อม ข้อมูลการพยากรณ์ปริมาณการผลิตไฟฟ้าจากแผงโซลาร์เซลล์  และการแสดงผลข้อมูลผ่านเว็บไซต์ (Web Site)
    \item[2.] นายปุณณวิชญ์ เดชอินทร์: รับผิดชอบในการออกแบบและพัฒนาสถานีวัดสภาพอากาศขนาดย่อมและส่งข้อมูลที่วัดได้ไปยังระบบจัดเก็บข้อมูล
    \item[3.] นายรัชพล แปรงทอง: รับผิดชอบในการออกแบบและพัฒนาระบบพยากรณ์ปริมาณการผลิตไฟฟ้าจากแผงโซลาร์เซลล์ล่วงหน้า 24 ชั่วโมง
\end{itemize}



\section{\ifenglish%
Impacts of this project on society, health, safety, legal, and cultural issues
\else%
ผลกระทบด้านสังคม สุขภาพ ความปลอดภัย กฎหมาย และวัฒนธรรม
\fi}
\begin{itemize}
    \item[] โครงงานระบบบริหารจัดการพลังงานอัจฉริยะมีประโยชน์ต่อสังคมในยุคที่ความต้องการพลังงานเพิ่มสูงขึ้นอย่างต่อเนื่อง โดยระบบสามารถนำข้อมูลสภาพอากาศและการผลิตพลังงานจากแผงโซลาร์เซลล์มาวิเคราะห์และพยากรณ์ เพื่อช่วยให้ผู้ใช้งานสามารถวางแผนการใช้พลังงานได้อย่างเหมาะสม ลดการพึ่งพาพลังงานจากภายนอก และส่งเสริมการใช้พลังงานอย่างมีประสิทธิภาพ นอกจากนี้ยังช่วยสร้างความตระหนักรู้เกี่ยวกับการอนุรักษ์พลังงานและการลดการปล่อยก๊าซเรือนกระจก ซึ่งส่งผลดีต่อสังคมโดยรวมในระยะยาว
    \item[] ในด้านสุขภาพ การนำพลังงานสะอาดมาใช้จากแผงโซลาร์เซลล์มีส่วนช่วยลดการปล่อยมลพิษทางอากาศจากแหล่งผลิตไฟฟ้าแบบดั้งเดิม ทำให้คุณภาพอากาศดีขึ้นและลดความเสี่ยงต่อโรคระบบทางเดินหายใจของประชาชน
    \item[] เมื่อพิจารณาในด้านความปลอดภัย ระบบมีการใช้อุปกรณ์ภาคสนาม เช่น เซนเซอร์ ที่มีมาตรฐานความปลอดภัยสูงและทนทานต่อสภาพอากาศที่มีความรุนแรง เพื่อป้องกันอุบัติเหตุที่อาจเกิดขึ้นจากการใช้งานอุปกรณ์
    \item[] ในด้านกฎหมาย การพัฒนาและใช้งานระบบมีความสอดคล้องกับข้อกำหนดและกฎหมายที่เกี่ยวข้อง เช่น กฎหมายด้านพลังงานไฟฟ้า รวมถึงมาตรฐานความปลอดภัยของอุปกรณ์ไฟฟ้าและอิเล็กทรอนิกส์ เพื่อให้การดำเนินงานเป็นไปอย่างถูกต้องและสามารถนำไปใช้งานได้จริงในเชิงปฏิบัติ
    \item[] สุดท้าย ในด้านวัฒนธรรม โครงงานนี้มีส่วนช่วยส่งเสริมค่านิยมด้านการใช้พลังงานอย่างยั่งยืนและความรับผิดชอบต่อสิ่งแวดล้อม ซึ่งสอดคล้องกับแนวทางการพัฒนาที่ยั่งยืนในสังคมไทย แม้ว่าการนำเทคโนโลยีใหม่มาใช้งานอาจต้องอาศัยการปรับตัวของผู้ใช้งาน แต่หากมีการให้ความรู้และสร้างความเข้าใจอย่างเหมาะสม จะช่วยให้เกิดการยอมรับและสามารถใช้งานระบบได้อย่างมีประสิทธิภาพ
\end{itemize}

